\label{sec:limitations_body}

\paragraph{Generalisability.}\label{ssec:lim_generalize}
Stack Overflow is one Q\&A platform in one domain (programming).
The decision point we identify---the search-for-help moment that
historically ended in a public query and now ends in an LLM
conversation---is a portable object that generates testable
hypotheses for medical, legal \citep{choi2022chatgpt}, and
academic Q\&A, implementable with the same design and API; we do
not undertake them here because each domain needs its own
substitutability taxonomy. We also do not observe other public
channels (Reddit, Discord, forums), so the DDD estimates changes
\emph{on} Stack Overflow but cannot distinguish movement to
private LLMs from movement to other public channels.

\paragraph{Input margin, not innovation outcomes.}\label{ssec:lim_scope}
The DDD estimates an innovation \emph{input}
(\S\ref{ssec:theory_qa_as_input}), not patents, products, entry,
or productivity; the informative horizon for downstream
propagation is probably 5--10 years and our window is two.
Heterogeneity by user tenure---the margin \citet{xue2026chatgpt}
study---is not separately observed in the aggregated cell-level
panel; a follow-up using the Stack Exchange data dump
(zero-budget but storage-intensive) could close this gap.

\paragraph{Measurement of substitutability.}\label{ssec:lim_measurement}
The binary classification is validated but imperfect (Fleiss
$\kappa=0.52$; seven-way $\kappa=0.40$;
\S\ref{sec:validation}), and the negative coefficient on
advanced-architecture (\S\ref{ssec:mech_qtype_anomaly}) suggests
the boundary moves as capabilities improve. The pre-treatment
proxy is not a direct LLM-capability benchmark; we report the
embedding-answerability treatment as a co-primary check, but a
direct pass-rate measure (a local open-weight model over sampled
pre-period questions) would be more precise and is a natural next
step within the zero-budget rule.

\paragraph{Pre-trends.}\label{ssec:lim_pretrend}
Exact parallel trends are rejected; we treat the pre-period
convergence, placebo sign-reversal, Copilot placebo, and the
HonestDiD $\Delta^{\text{RM}}$ profile (robust CI first contains
zero at $\bar M=0.75$; \S\ref{ssec:id_honest_did}) as evidence
against a simple secular-trend explanation, not proof of parallel
trends, and describe the causal reading as moderate rather than
strong. A quadratic trend control (in the replication package)
absorbs more of the effect mechanically by mimicking the shock
shape; we report it as a limit case.

\paragraph{No welfare conclusion.}\label{ssec:lim_welfare}
We bound one public-cost channel from above. The private
productivity gains documented elsewhere are first-order and
plausibly an order of magnitude larger
(\S\ref{ssec:pol_welfare}), and the long-run model-degradation
implications \citep{shumailov2024curse} are unresolved.

\paragraph{Scope of the artefact.}\label{ssec:lim_artefact}
The early-warning system is a reproducible scoring-and-alerting
script (\S\ref{ssec:pol_monitor}); a production dashboard and a
usability evaluation with platform moderators would be required
to claim a deployed, human-evaluated system, and neither was
available for this study. (The 2025 panel extension earlier
flagged as future work is now complete for all 100 tags and used
in \S\ref{ssec:id_threats}.)
